% TOP_PROBLEM: {"id": 30005276, "title": "Density Hypothesis for Riemann-Zeta Zeros", "category_id": 1, "category": {"id": 1, "name": "number_theory", "display_name": "Number Theory", "description": "Properties of integers, prime numbers, Diophantine equations.", "slug": "number-theory", "order_index": 1, "created_at": "2026-07-31T15:26:25.670Z"}, "status": "partially_solved"}
% FIRST_POSTED: 2026-10-06
% ATTEMPT_STATUS: unresolved
% OPENAI_CLAIM: {"schema_version":1,"source_repo":"https://github.com/openai/math","source_commit":"adc7f1241b42e322a6451854ab7e4b4c146bf78a","release_date":"2026-10-06","match":"partial","resolution":"partial","independently_reviewed":false,"second_pass":"n/a","families":[{"family":"003","title":"The quasi-Riemann hypothesis","subject":"Number theory","manuscripts":[{"title":"The Quasi-Riemann Hypothesis","pdf_path":"preprints/The-Quasi-Riemann-Hypothesis-October-5-2026/paper2.pdf","readme_path":"preprints/The-Quasi-Riemann-Hypothesis-October-5-2026/README.md","date":"2026-10-05","theorem_ref":"thm:main (paper2.tex:73)","statement_sources":[{"path":"preprints/The-Quasi-Riemann-Hypothesis-October-5-2026/build/paper2.tex","line":73,"label":"thm:main"}]},{"title":"The Quasi-Riemann Hypothesis: A Zero-Free Half-Plane Re (s)>7/8","pdf_path":"preprints/The-Quasi-Riemann-Hypothesis-September-30-2026/paper.pdf","readme_path":"preprints/The-Quasi-Riemann-Hypothesis-September-30-2026/README.md","date":"2026-09-30","theorem_ref":"thm:main (paper.tex:106)","statement_sources":[{"path":"preprints/The-Quasi-Riemann-Hypothesis-September-30-2026/build/paper.tex","line":106,"label":"thm:main"}]}],"lean":{"doc_path":"lean/docs/003.md","comparators":["lean/ComparatorChallenges/DirichletSevenEighths.json","lean/ComparatorChallenges/DirichletSevenEighths.lean","lean/ComparatorChallenges/HeckeSevenEighths.json","lean/ComparatorChallenges/HeckeSevenEighths.lean","lean/ComparatorChallenges/QuasiRiemannHypothesis.json","lean/ComparatorChallenges/QuasiRiemannHypothesis.lean","lean/ComparatorChallenges/SiegelZeros.json","lean/ComparatorChallenges/SiegelZeros.lean"],"declarations":["OAI.DirichletCharacter.LFunction_ne_zero_of_seven_eighths_lt_re","OAI.SevenEighths.HeckeFamily.LFunction_ne_zero_of_seven_eighths_lt_re","OAI.SiegelZeros.WeightedTorusJets.exists_absolute_real_zero_gap","OAI.riemannZeta_ne_zero_of_seven_eighths_lt_re"],"covers_main_theorem":false},"reasoning_trace":null}]}
\documentclass[11pt]{article}
\usepackage[margin=25mm]{geometry}
\usepackage{fontspec}
\setmainfont{Latin Modern Roman}
\usepackage{amsmath,amssymb,mathtools}
\usepackage{unicode-math}
\setmathfont{Latin Modern Math}
\usepackage{enumitem}
\usepackage{xurl,hyperref,bookmark}
\hypersetup{colorlinks=true,urlcolor=blue}
\setcounter{secnumdepth}{0}
\setlength{\parindent}{0pt}
\setlength{\parskip}{5pt}
\setlength{\emergencystretch}{3em}
\begin{document}
\section*{Density Hypothesis for Riemann-Zeta Zeros}

\subsection{OpenAI's claim}
OpenAI presents the following claims relevant to this question, with manuscript numbers following paper order in Sources. In manuscript 1 (family 003; 2026-10-05), the asserted result is that zeta, Dirichlet functions, and finite-order Eisenstein-field Hecke families admit a zero-free region with boundary $11/12$; In manuscript 2 (family 003; 2026-09-30), the asserted result is that zeta, Dirichlet $L$-functions, and finite-order Hecke $L$-functions for $\mathbb Q(\sqrt{-3})$ have no zeros with real part greater than $7/8$, with poles excluded.

\subsection{Mathematical statement of the claimed result}
\paragraph{Manuscript 1.}
For $K=\mathbb Q(\sqrt{-3})$ and every finite-order Hecke character $\chi$ over $K$, the claimed zero-free assertion, with meromorphic poles excluded from evaluation, is
\[
 \operatorname{Re}s>11/12\quad\Longrightarrow\quad L_K(s,\chi)\ne0.
\]
It includes every Dirichlet $L$-function and the Riemann zeta function. The accompanying prime-progression claim uses $\pi(x;q,a)$ and asserts, with an absolute effective constant,
\[
 \sup_{\substack{1\le q\le x\\(a,q)=1}}
 \left|\pi(x;q,a)-\frac1{\varphi(q)}\int_2^x\frac{dt}{\log t}\right|
 \ll x^{11/12}\log x\quad(x\ge2).
\]
The critical-line requirement $\operatorname{Re}\rho=1/2$ is not a consequence of this half-plane claim.

\paragraph{Manuscript 2.}
For every finite-order Hecke character of $F=\mathbb Q(\sqrt{-3})$, every Dirichlet character and the zeta function, the principal claim is
\[
 \operatorname{Re}s>7/8\quad\Longrightarrow\quad L(s)\ne0,
\]
where a principal character's pole at $s=1$ is allowed. The weaker displayed $11/12$ theorem follows within this scope. The cited arithmetic consequence gives absolute $A,C>0$ such that every odd prime $p$ has least positive quadratic nonresidue
\[
 n(p)\le C(\log p)^A.
\]
Thus $n(p)\ll_{\delta} p^{\delta}$ for each fixed $\delta>0$. Given binary $p$ and $a\in\mathbb F_p$, a deterministic algorithm of time polynomial in $\log p$ is further claimed to return a square root of $a$ or correctly declare it nonexistent. The half-plane assertion leaves zeros between the critical line and $7/8$ untreated.

\subsection{Outline of the argument}
For manuscript 1 (family 003; 2026-10-05), twisted M\"obius sums are enlarged to a family before Poisson summation. A recursive estimate for the dual mean square, followed by inversion of cube factors, supplies the saving used to exclude a reciprocal-function singularity.

For manuscript 2 (family 003; 2026-09-30), smoothed character sums are estimated in two ways: directly by large-sieve and recursive moment bounds, and through a Mellin transform involving the reciprocal of an $L$-function. Poisson summation and sextic residue characters align the two expressions so that their savings force continuation of the reciprocal.

\subsection{Scope relative to this problem}
The requested density exponent must hold across the full range starting at real part $1/2$. A zero-free half-plane makes the count vanish only above its fixed boundary, so it supplies no density estimate in the rest of that range.

\subsection{Formal verification}
For manuscript 2 (family 003; 2026-09-30), OpenAI lists Lean statements for this assertion: zeta, Dirichlet $L$-functions, and finite-order Hecke $L$-functions for $\mathbb Q(\sqrt{-3})$ have no zeros with real part greater than $7/8$, with poles excluded.

This coverage does not establish a formal proof of the complete problem as stated here.

No Lean formalization is listed at this snapshot for manuscript 1 (family 003; 2026-10-05).

The linked family Comparator checks are \texttt{DirichletSevenEighths.lean}, \texttt{HeckeSevenEighths.lean}, \texttt{QuasiRiemannHypothesis.lean}, \texttt{SiegelZeros.lean}.

Their named declarations include \texttt{OAI.DirichletCharacter.LFunction\_ne\_zero\_of\_seven\_eighths\_lt\_re}, \texttt{OAI.SevenEighths.HeckeFamily.LFunction\_ne\_zero\_of\_seven\_eighths\_lt\_re}, \texttt{OAI.SiegelZeros.WeightedTorusJets.dirichletRealZeroBound\_proof}, \texttt{OAI.SiegelZeros.WeightedTorusJets.exists\_absolute\_real\_zero\_gap}, \texttt{OAI.riemannZeta\_ne\_zero\_of\_seven\_eighths\_lt\_re}.

This site has not executed the Lean code. Described statement coverage is not an independent proof check.

\subsection{Status caveat}
The manuscripts were produced by an internal OpenAI model. OpenAI cautions that results without formalization can contain mistakes. This site has not independently checked the proofs; the label records the scope of OpenAI's claim.

\subsection{Sources}
\begin{itemize}
\item OpenAI, \emph{The Quasi-Riemann Hypothesis}, 2026-10-05; citation key \texttt{OAI:The-Quasi-Riemann-Hypothesis-October-5-2026}. \href{https://github.com/openai/math/blob/adc7f1241b42e322a6451854ab7e4b4c146bf78a/preprints/The-Quasi-Riemann-Hypothesis-October-5-2026/paper2.pdf}{PDF (GitHub)}; \href{https://github.com/openai/math/raw/adc7f1241b42e322a6451854ab7e4b4c146bf78a/preprints/The-Quasi-Riemann-Hypothesis-October-5-2026/paper2.pdf}{PDF (direct)}; \href{https://github.com/openai/math/blob/adc7f1241b42e322a6451854ab7e4b4c146bf78a/preprints/The-Quasi-Riemann-Hypothesis-October-5-2026/README.md}{Manuscript page and citation}; \href{https://github.com/openai/math/blob/main/preprints/The-Quasi-Riemann-Hypothesis-October-5-2026/paper2.pdf}{Latest version}; Statement source: \href{https://github.com/openai/math/blob/adc7f1241b42e322a6451854ab7e4b4c146bf78a/preprints/The-Quasi-Riemann-Hypothesis-October-5-2026/build/paper2.tex\#L73}{thm:main}.
\item OpenAI, \emph{The Quasi-Riemann Hypothesis: A Zero-Free Half-Plane Re (s)>7/8}, 2026-09-30; citation key \texttt{OAI:The-Quasi-Riemann-Hypothesis-September-30-2026}. \href{https://github.com/openai/math/blob/adc7f1241b42e322a6451854ab7e4b4c146bf78a/preprints/The-Quasi-Riemann-Hypothesis-September-30-2026/paper.pdf}{PDF (GitHub)}; \href{https://github.com/openai/math/raw/adc7f1241b42e322a6451854ab7e4b4c146bf78a/preprints/The-Quasi-Riemann-Hypothesis-September-30-2026/paper.pdf}{PDF (direct)}; \href{https://github.com/openai/math/blob/adc7f1241b42e322a6451854ab7e4b4c146bf78a/preprints/The-Quasi-Riemann-Hypothesis-September-30-2026/README.md}{Manuscript page and citation}; \href{https://github.com/openai/math/blob/main/preprints/The-Quasi-Riemann-Hypothesis-September-30-2026/paper.pdf}{Latest version}; Statement source: \href{https://github.com/openai/math/blob/adc7f1241b42e322a6451854ab7e4b4c146bf78a/preprints/The-Quasi-Riemann-Hypothesis-September-30-2026/build/paper.tex\#L106}{thm:main}.
\item Family 003: \href{https://github.com/openai/math/blob/adc7f1241b42e322a6451854ab7e4b4c146bf78a/lean/docs/003.md}{Lean scope}.
\item \href{https://github.com/openai/math/blob/adc7f1241b42e322a6451854ab7e4b4c146bf78a/lean/ComparatorChallenges/DirichletSevenEighths.json}{Comparator DirichletSevenEighths.json}.
\item \href{https://github.com/openai/math/blob/adc7f1241b42e322a6451854ab7e4b4c146bf78a/lean/ComparatorChallenges/DirichletSevenEighths.lean}{Comparator DirichletSevenEighths.lean}.
\item \href{https://github.com/openai/math/blob/adc7f1241b42e322a6451854ab7e4b4c146bf78a/lean/ComparatorChallenges/HeckeSevenEighths.json}{Comparator HeckeSevenEighths.json}.
\item \href{https://github.com/openai/math/blob/adc7f1241b42e322a6451854ab7e4b4c146bf78a/lean/ComparatorChallenges/HeckeSevenEighths.lean}{Comparator HeckeSevenEighths.lean}.
\item \href{https://github.com/openai/math/blob/adc7f1241b42e322a6451854ab7e4b4c146bf78a/lean/ComparatorChallenges/QuasiRiemannHypothesis.json}{Comparator QuasiRiemannHypothesis.json}.
\item \href{https://github.com/openai/math/blob/adc7f1241b42e322a6451854ab7e4b4c146bf78a/lean/ComparatorChallenges/QuasiRiemannHypothesis.lean}{Comparator QuasiRiemannHypothesis.lean}.
\item \href{https://github.com/openai/math/blob/adc7f1241b42e322a6451854ab7e4b4c146bf78a/lean/ComparatorChallenges/SiegelZeros.json}{Comparator SiegelZeros.json}.
\item \href{https://github.com/openai/math/blob/adc7f1241b42e322a6451854ab7e4b4c146bf78a/lean/ComparatorChallenges/SiegelZeros.lean}{Comparator SiegelZeros.lean}.
\item \href{https://github.com/openai/math/blob/adc7f1241b42e322a6451854ab7e4b4c146bf78a/CONTENTS.md}{OpenAI catalogue entry}.
\end{itemize}
\end{document}
